\documentclass[10pt,a4paper]{amsart}
\usepackage[margin=19mm]{geometry}
\usepackage{amsmath,amssymb,mathtools}
\usepackage[scr=rsfs,cal=euler]{mathalfa}
\usepackage{xfrac}
\usepackage[unicode,hidelinks,bookmarks=false]{hyperref}
\newcommand{\ie}{i.e.\ }
\newcommand{\cf}{cf.\ }
\newcommand{\resp}{respectively\ }
\newcommand{\Subsection}{\subsection}
\providecommand{\nobreakdash}{\nobreak\hbox{-}}
\setlength{\emergencystretch}{2em}
\multlinegap0pt
\usepackage{siunitx}
\usepackage{siunitx}
\usepackage{siunitx}
\usepackage{siunitx}
\usepackage{amssymb}
\usepackage{mathtools}
\usepackage{bm}
\usepackage{tensor}
\usepackage{dsfont}
\usepackage{xcolor}
\usepackage{siunitx}
\usepackage{tikz}
\usepackage[shortlabels]{enumitem}
\usepackage{braket}
\usepackage{csquotes}
\newtheorem{theorem}{Theorem}
\DeclareMathOperator{\Sing}{Sing}
\newcommand{\cF}{\mathcal{F}}
\title{Excess logarithmic residues for foliations by curves and applications}
\author{Alana Cavalcante, Maur\'icio Corr\^ea, Fernando Louren\c co, Elaheh Shahsavaripour}
\date{Source: arXiv:2604.25395v1 (CC BY 4.0)}
\begin{document}
\maketitle

\noindent\emph{Source and licence.} Excess logarithmic residues for foliations by curves and applications. Version arXiv:2604.25395v1 (CC BY 4.0), distributed by arXiv under the Creative Commons Attribution 4.0 International licence, http://creativecommons.org/licenses/by/4.0/, as recorded in the arXiv licence metadata for this exact version and shown on the abstract page https://arxiv.org/abs/2604.25395v1. Full licence notice: \texttt{licenses/arxiv-2604.25395-CC-BY-4.0.txt}.\par\smallskip

\noindent\textit{Editorial scope.} Counted unit: the article's theorem thm:residues (source0.tex lines 1383--1430). The article's complete original proof of it is delivered with it (source0.tex lines 1432--1546, one \texttt{proof} environment of 115 lines). No mathematics is written, repaired or re-derived: the statement and the proof are the article's own lines, in the article's own notation, and no retained line is substituted or re-flowed.  No further source-local context is needed: every cross-reference of every retained line resolves inside the two delivered spans.  Nothing was omitted from any retained span, so the package carries no omission record.  No retained line contains a citation, so the wrapper carries no bibliography and no bibliography entry.  No retained line contains a figure environment, an image-inclusion command or drawing code, so no image is delivered and the package declares no asset dependency.  Editorial formatting only, in addition: an AMS \texttt{amsart} preamble that restates the article's own declarations and macros used by the retained lines, namely the environments \texttt{theorem} on separate counters, together with the standard packages those lines need.  The retained environments print in this document as Theorem 1 in order of appearance, whereas the article prints the same items with the numbering of its own full document.  The complete native source of the article as distributed in the arXiv e-print for this exact version is bundled unchanged as \texttt{sources/199301/source0.tex}.
\begin{theorem}\label{thm:residues}
Let \(\cF\) be a one-dimensional holomorphic foliation on \(X\) with only
isolated singularities, logarithmic along \(D\). Let \(Z_D(\cF)\) denote the
finite set of isolated zeroes of the vector field induced by \(\cF\) on
\(D_{\mathrm{reg}}\). When these zeroes coincide with
\(\Sing(\cF)\cap D_{\mathrm{reg}}\), we write the sums simply over
\(\Sing(\cF)\cap D\).
For \(i=0\), one has the logarithmic formula
\[
\sum_{p\in \Sing(\cF)}
\widehat{\operatorname{Res}}^{\log}_{c_1^n}(\cF,D,p)
=
\int_X c_1\bigl(N^{\log}_{\cF/X}\bigr)^n,
\]
and the variational identity
\[
\sum_{p\in \Sing(\cF)}
\operatorname{Var}_{c_1^n}(\cF,D,p)
=
\int_X
\Bigl[
c_1(N_{\cF})^n
-
c_1\bigl(N^{\log}_{\cF/X}\bigr)^n
\Bigr].
\]
For every \(1\leq i\leq n-1\), one has
\[
\sum_{p\in Z_D(\cF)}
\widehat{\operatorname{Res}}^{\log}_{c_1^{n-i},\,D^i}
(\cF,D,p)
=
\int_X c_1\bigl(N^{\log}_{\cF/X}\bigr)^{n-i}D^i,
\]
and
\[
\sum_{p\in Z_D(\cF)}
\operatorname{Var}_{c_1^{n-i},\,D^i}
(\cF,D,p)
=
\int_X
\Bigl[
c_1(N_{\cF})^{n-i}
-
c_1\bigl(N^{\log}_{\cF/X}\bigr)^{n-i}
\Bigr]D^i.
\]
\end{theorem}

\begin{proof}
We prove the logarithmic identities. The ordinary identities are obtained in
the same way by replacing the logarithmic normal sheaf by the ordinary normal
sheaf, and the variational identities follow by subtraction.
First consider \(i=0\). Let
\[
V_0=\bigcup_{p\in \Sing(\cF)}B_p
\]
be a union of pairwise disjoint sufficiently small balls centred at the isolated
singularities of \(\cF\), and put \(L_p=\partial B_p\). On the punctured ball
\(B_p^\ast=B_p\setminus\{p\}\), choose Pfaff data as in the preceding
construction:
\[
\omega_{12}=\eta_1\wedge\cdots\wedge\eta_{n-1},
\qquad
\omega_1=\omega_2\wedge\omega_{12}
=
h\,\eta_1\wedge\cdots\wedge\eta_{n-1},
\]
with \(\omega_1\) defining \(\cF\). The integrability relations give the trace
forms \(\theta^{12}\) and \(\widetilde\theta^{12}\), and we set
\[
F_{12}^{\log}:=d\widetilde\theta^{12}.
\]
The form
\[
\psi_0^{\log}
:=
(2\pi i)^{-n}
\widetilde\theta^{12}\wedge
\bigl(F_{12}^{\log}\bigr)^{n-1}
\]
satisfies
\[
d\psi_0^{\log}
=
(2\pi i)^{-n}
\bigl(F_{12}^{\log}\bigr)^n,
\]
which represents \(c_1(N^{\log}_{\cF/X})^n\). Since \(\cF\) has dimension one,
its codimension is \(n-1\), and Bott's vanishing theorem localizes this
degree-\(n\) characteristic class on \(\Sing(\cF)\). Therefore Stokes' theorem
on \(X\setminus V_0\), followed by the limit as the radii of the balls tend to
zero, gives
\[
\sum_{p\in \Sing(\cF)}
\widehat{\operatorname{Res}}^{\log}_{c_1^n}(\cF,D,p)
=
\sum_{p\in \Sing(\cF)}
\int_{L_p}\psi_0^{\log}
=
\int_X c_1\bigl(N^{\log}_{\cF/X}\bigr)^n.
\]

Now let \(1\leq i\leq n-1\). In this case the factor \(D^i\) makes the
localization relative to the divisor. On \(D_{\mathrm{reg}}\), the logarithmic
vector field induces a one-dimensional foliation of codimension \(n-2\). The
relative characteristic polynomial has degree \(n-1\), and Bott's vanishing
theorem on \(D_{\mathrm{reg}}\) localizes the corresponding relative class at
the isolated zeroes of the induced vector field, namely at \(Z_D(\cF)\).
Let
\[
V_D=\bigcup_{p\in Z_D(\cF)}B_p
\]
be a union of pairwise disjoint sufficiently small balls centred at these
points. On each punctured ball choose the same local Pfaff data as above and
also a local primitive \(\beta\) of the fixed closed representative \(F_D\) of
the class \(D\), so that \(d\beta=F_D\). Set
\[
\psi_i^{\log}
:=
(2\pi i)^{-n}
\widetilde\theta^{12}\wedge
F_D^i\wedge
\bigl(F_{12}^{\log}\bigr)^{n-i-1}.
\]
Since \(dF_D=dF_{12}^{\log}=0\), one has
\[
d\psi_i^{\log}
=
(2\pi i)^{-n}
F_D^i\wedge
\bigl(F_{12}^{\log}\bigr)^{n-i}.
\]
This form represents
\[
c_1\bigl(N_{\cF/X}^{\log}\bigr)^{n-i}D^i.
\]
By the relative Bott localization described above, this class is represented
by a current supported on \(\overline V_D\). Applying Stokes' theorem on
\(X\setminus V_D\) and then letting the radii of the balls tend to zero, we get
\[
\sum_{p\in Z_D(\cF)}
\widehat{\operatorname{Res}}^{\log}_{c_1^{n-i},\,D^i}(\cF,D,p)
=
\sum_{p\in Z_D(\cF)}
\int_{L_p}\psi_i^{\log}
=
\int_X
c_1\bigl(N_{\cF/X}^{\log}\bigr)^{n-i}D^i.
\]
The proof of the ordinary identities is identical, replacing
\(\widetilde\theta^{12}\) and \(F_{12}^{\log}\) by the corresponding ordinary
trace form and curvature representative of \(c_1(N_{\cF})\). Finally, by
definition,
\[
\operatorname{Var}_{c_1^{n-i},D^i}(\cF,D,p)
=
\operatorname{Res}_{c_1^{n-i},D^i}(\cF,D,p)
-
\widehat{\operatorname{Res}}^{\log}_{c_1^{n-i},D^i}(\cF,D,p).
\]
Thus the variational identities follow by subtracting the logarithmic
identities from the ordinary ones.
\end{proof}

\end{document}
